% Einheit 2 von 4 · Normalform und p-q-Formel (Katalog Einheit 3)
\hypertarget{e2}{}\einheitenkopf{Einheit 2 von 4 · Normalform und p-q-Formel}
\zweigzeile{Hier lernst du, eine quadratische Gleichung zu ordnen, zu normieren und mit der p-q-Formel zu lösen · neu in diesem Jahr, je nach Buch erst Klasse 10 · P10 oft · baut auf: Wurzelziehen und Lösbarkeit (Einheit 1), Terme ordnen, Ausmultiplizieren und binomische Formeln}
\renewcommand{\theaufgabe}{\arabic{aufgabe}}\setcounter{aufgabe}{10}

\begin{aufgabe}{Ich erkenne an der Form einer Gleichung, welcher Lösungsweg passt. Kreuze an. Du musst nichts lösen.}
\begin{geruest}
\gz{a}{$x^2 - 5x + 4 = 0$}{\kreuz{Wurzelziehen}\kreuz{Rückwärtsrechnen}\kreuz{p-q-Formel}}
\gz{b}{$x^2 = 13$}{\kreuz{Wurzelziehen}\kreuz{Rückwärtsrechnen}\kreuz{p-q-Formel}}
\gz{c}{$(x + 2)^2 = 11$}{\kreuz{Wurzelziehen}\kreuz{Rückwärtsrechnen}\kreuz{p-q-Formel}}
\gz{d}{$3x^2 - 12 = 0$}{\kreuz{Wurzelziehen}\kreuz{Rückwärtsrechnen}\kreuz{p-q-Formel}}
\gz{e}{$2x^2 + 3x - 1 = 0$}{\kreuz{Wurzelziehen}\kreuz{Rückwärtsrechnen}\kreuz{p-q-Formel}}
\end{geruest}
\end{aufgabe}

\begin{aufgabe}{Ich kann in der Normalform $x^2 + px + q = 0$ die Zahlen $p$ und $q$ mit Vorzeichen ablesen. Kreise $p$ und $q$ ein und schreibe sie auf. Du musst nichts ausrechnen.}
\begin{geruest}
\gz{a}{$x^2 + 7x + 10 = 0$}{\feld{p}\feld{q}}
\gz{b}{$x^2 - 3x + 2 = 0$}{\feld{p}\feld{q}}
\gz{c}{$x^2 + 4x - 7 = 0$}{\feld{p}\feld{q}}
\gz{d}{$x^2 - x - 12 = 0$}{\feld{p}\feld{q}}
\gz{e}{$x^2 - 11x = 0$}{\feld{p}\feld{q}}
\end{geruest}
\end{aufgabe}

\begin{aufgabe}{Ich erkenne, ob die Zahl vor $x^2$ eins ist oder ob ich vorher teilen muss. Kreuze an. Musst du teilen, schreibe die Zahl auf, durch die du teilst. Du musst nichts rechnen.}
\begin{geruest}
\gz{a}{$x^2 + 3x - 4 = 0$}{\kreuz{Zahl vor $x^2$ ist eins}\kreuz{teilen durch} \leerfeld}
\gz{b}{$2x^2 + 6x - 8 = 0$}{\kreuz{Zahl vor $x^2$ ist eins}\kreuz{teilen durch} \leerfeld}
\gz{c}{$5x^2 - 10x + 5 = 0$}{\kreuz{Zahl vor $x^2$ ist eins}\kreuz{teilen durch} \leerfeld}
\gz{d}{$-x^2 + 4x + 5 = 0$}{\kreuz{Zahl vor $x^2$ ist eins}\kreuz{teilen durch} \leerfeld}
\gz{e}{$0{,}5x^2 + x - 3 = 0$}{\kreuz{Zahl vor $x^2$ ist eins}\kreuz{teilen durch} \leerfeld}
\end{geruest}
\end{aufgabe}

\begin{aufgabe}{Ich kann eine Gleichung in Normalform mit der p-q-Formel lösen (kommt je nach Buch erst nächstes Jahr). Die Formel lautet $x_{1,2} = -\frac{p}{2} \pm \sqrt{\left(\frac{p}{2}\right)^2 - q}$. Lies $p$ und $q$ ab, setze ein und gib beide Lösungen an.}
\beispiel{x^2 + 2x - 8 &= 0 && p = 2,\ q = -8 \\ x_{1,2} &= -1 \pm \sqrt{1^2 + 8} \\ x_{1,2} &= -1 \pm 3 \\ x_1 &= 2, \quad x_2 = -4}
\begin{gleichungsraster}[3]
\gl{x^2 + 6x + 8 = 0} & \gl{x^2 + 10x + 9 = 0} \\ \rule{0pt}{2mm} & \\
\gl{x^2 + 6x + 5 = 0} & \gl{x^2 + 10x + 16 = 0} \\ \rule{0pt}{2mm} & \\
\gl{x^2 - 6x + 8 = 0} & \gl{x^2 + 2x - 24 = 0} \\ \rule{0pt}{2mm} & \\
\end{gleichungsraster}
\end{aufgabe}

\begin{aufgabe}{Ich kann eine Gleichung erst in Normalform bringen und dann mit der p-q-Formel lösen. Bringe alles auf eine Seite, ordne nach $x^2$, $x$ und Zahl und teile durch die Zahl vor $x^2$, wenn sie nicht eins ist.}
\beispiel{2x^2 + 12x &= 32 &&\mid -32 \\ 2x^2 + 12x - 32 &= 0 &&\mid :2 \\ x^2 + 6x - 16 &= 0 && p = 6,\ q = -16 \\ x_{1,2} &= -3 \pm \sqrt{9 + 16} = -3 \pm 5 \\ x_1 &= 2, \quad x_2 = -8}
\begin{gleichungsraster}[4]
\gl{x^2 - 4x = 5} & \gl{x^2 + 5x = x + 12} \\ \rule{0pt}{2mm} & \\
\gl{5x^2 - 20x + 15 = 0} & \gl{-x^2 + 8x - 12 = 0} \\ \rule{0pt}{2mm} & \\
\end{gleichungsraster}
\end{aufgabe}

\begin{aufgabe}{Ich kann eine Gleichung in Normalform mit der p-q-Formel lösen – weiter. Löse mit der p-q-Formel.}
\begin{gleichungsraster}[3]
\gl{x^2 + 3x - 10 = 0} & \gl{x^2 - 10x + 25 = 0} \\ \rule{0pt}{2mm} & \\
\gl{x^2 + 4x + 7 = 0} & \\ \rule{0pt}{2mm} & \\
\end{gleichungsraster}
Berechne nur den Wert unter der Wurzel. Man nennt ihn Diskriminante. Gib an, wie viele Lösungen die Gleichung hat.
\begin{geruest}
\gz{d}{$x^2 + 5x + 3 = 0$}{Wert unter der Wurzel: \leerfeld \quad Anzahl: \leerfeld}
\end{geruest}
Löse. Runde die Lösungen auf zwei Stellen nach dem Komma.
\begin{gleichungsraster}[3]
\gl{x^2 - 4x + 1 = 0} & \\ \rule{0pt}{2mm} & \\
\end{gleichungsraster}
Löse und mache mit einer Lösung die Probe.
\begin{gleichungsraster}[4]
\gl{x^2 + 2x - 35 = 0} & \\ \rule{0pt}{2mm} & \\
\end{gleichungsraster}
\end{aufgabe}

\begin{aufgabe}{Ich kann Klammern zuerst auflösen und den passenden Lösungsweg wählen. Löse die Klammer auf, bringe die Gleichung in Normalform und löse sie.}
\begin{gleichungsraster}[4]
\gl{(x + 1)^2 = 3x + 7} & \gl{x \cdot (x + 4) = 21} \\ \rule{0pt}{2mm} & \\
\end{gleichungsraster}
Wähle einen kurzen Lösungsweg: Wurzelziehen, Rückwärtsrechnen oder p-q-Formel. Löse.
\begin{gleichungsraster}[3]
\gl{2x^2 - 50 = 0} & \gl{(x - 1)^2 = 16} \\ \rule{0pt}{2mm} & \\
\end{gleichungsraster}
Berechne die Nullstellen der Funktion $f$. (P10 2020)
\begin{gleichungsraster}[4]
\gl{f(x) = 2x^2 - 12x + 10} & \\ \rule{0pt}{2mm} & \\
\end{gleichungsraster}
\end{aufgabe}

\begin{aufgabe}{Ich kann die Schnittpunkte einer Parabel mit einer Geraden berechnen. Die Parabel gehört zu $y = x^2 - 5$, die Gerade zu $y = -x + 1$.}

\begin{ksys}[xmin=-4,xmax=4,ymin=-6,ymax=6,ablesen]
\parabel{1}{0}{-5}{}
\gerade{-1}{1}{}
\end{ksys}

\begin{teile}
\teil Lies die Schnittpunkte $A$ und $B$ ab.\\ $A($\leerfeld$|$\leerfeld$)$ \quad $B($\leerfeld$|$\leerfeld$)$
\end{teile}
Berechne die Schnittpunkte der Parabel und der Geraden: Setze die beiden Terme gleich, bringe die Gleichung in Normalform, löse sie und berechne die $y$-Werte.
\begin{gleichungsraster}[4]
\gl{x^2 - 5 = -x + 1} & \\ \rule{0pt}{2mm} & \\
\end{gleichungsraster}
Gegeben sind die Funktionen $g$ und $k$. Berechne die Koordinaten der Schnittpunkte ihrer Graphen. (P10 2022)
\begin{gleichungsraster}[4]
\gl{g(x) = (x - 1)^2 - 3, \quad k(x) = x + 2} & \\ \rule{0pt}{2mm} & \\
\end{gleichungsraster}
\end{aufgabe}

\begin{aufgabe}{Ich finde den Fehler beim Lösen mit der p-q-Formel.}
\begin{teile}
\teil Ali löst $2x^2 - 10x + 12 = 0$ so:
\rechnung{2x^2 - 10x + 12 &= 0 && p = -10,\ q = 12 \\ x_{1,2} &= 5 \pm \sqrt{25 - 12} \\ x_{1,2} &= 5 \pm \sqrt{13}}
Finde den Fehler, benenne ihn und korrigiere die Rechnung.
\end{teile}
Löse selbst.
\begin{gleichungsraster}[4]
\gl{5x^2 + 10x - 15 = 0} & \\ \rule{0pt}{2mm} & \\
\end{gleichungsraster}
\begin{teile}
\teil Sara löst $x^2 - 10x + 16 = 0$ so:
\rechnung{x^2 - 10x + 16 &= 0 && p = -10,\ q = 16 \\ x_{1,2} &= -5 \pm \sqrt{25 - 16} \\ x_{1,2} &= -5 \pm 3 \\ x_1 &= -2, \quad x_2 = -8}
Finde den Fehler, benenne ihn und korrigiere die Rechnung.
\end{teile}
Löse selbst.
\begin{gleichungsraster}[3]
\gl{x^2 - 14x + 40 = 0} & \\ \rule{0pt}{2mm} & \\
\end{gleichungsraster}
\end{aufgabe}

\begin{aufgabe}{Ich kann begründen, wann die p-q-Formel passt und wann es keine Lösung gibt.}
\begin{teile}
\teil Begründe, warum $x^2 + 2x + 5 = 0$ keine Lösung hat.
\teil Tim setzt bei $4x^2 + 8x - 4 = 0$ direkt $p = 8$ und $q = -4$ in die p-q-Formel ein. Begründe, warum das falsch ist.
\teil Entscheide, ohne die Lösungen auszurechnen, ob $x^2 - 6x + 9 = 0$ eine oder zwei Lösungen hat. Begründe.
\end{teile}
\end{aufgabe}

\begin{aufgabe}{Ich kann eine Sachfrage mit der p-q-Formel beantworten. Für den Anhalteweg $s$ (in m) eines Autos bei der Geschwindigkeit $v$ (in km/h) gilt die Faustformel $s = 0{,}3 \cdot v + 0{,}01 \cdot v^2$. Nach einem Unfall in einer Tempo-30-Zone misst die Polizei einen Anhalteweg von $40\,$m.}
\begin{teile}
\teil Stelle die Gleichung auf und bringe sie in Normalform.
\teil Berechne $v$.
\teil Begründe, warum nur eine der beiden Lösungen passt. War das Auto zu schnell?
\end{teile}
\end{aufgabe}
